\documentclass[12pt,a4paper]{article}
\usepackage[utf8]{inputenc}
\usepackage[spanish]{babel}
\usepackage[T1]{fontenc}
\usepackage{geometry}
\usepackage{graphicx}
\usepackage{xcolor}
\usepackage{titlesec}
\usepackage{enumitem}
\usepackage{fancyhdr}
\usepackage{tcolorbox}
\usepackage{fontawesome5}
\usepackage{hyperref}
\usepackage{booktabs}
\usepackage{array}
\usepackage{mdframed}

\geometry{top=2.5cm, bottom=2.5cm, left=2.5cm, right=2.5cm}

% ── Colores corporativos ──────────────────────────────────────────
\definecolor{primario}{HTML}{C9A96E}
\definecolor{oscuro}{HTML}{1C1C1E}
\definecolor{gris}{HTML}{6B6B6B}
\definecolor{fondo}{HTML}{FAF7F2}
\definecolor{alerta}{HTML}{E74C3C}
\definecolor{exito}{HTML}{2ECC71}
\definecolor{info}{HTML}{3498DB}
\definecolor{admin}{HTML}{8E44AD}

% ── Secciones ────────────────────────────────────────────────────
\titleformat{\section}{\Large\bfseries\color{oscuro}}{}{0em}{}[\color{primario}\titlerule]
\titleformat{\subsection}{\large\bfseries\color{oscuro}}{}{0em}{\faAngleRight\ }
\titleformat{\subsubsection}{\normalsize\bfseries\color{gris}}{}{0em}{}

% ── Cajas ────────────────────────────────────────────────────────
\tcbuselibrary{skins,breakable}
\newtcolorbox{notabox}{
  colback=fondo, colframe=primario, coltitle=white, fonttitle=\bfseries,
  title={\faInfoCircle\ Nota}, breakable, arc=4pt
}
\newtcolorbox{alertabox}{
  colback=red!5, colframe=alerta, coltitle=white, fonttitle=\bfseries,
  title={\faExclamationTriangle\ Importante}, breakable, arc=4pt
}
\newtcolorbox{pasobox}[1]{
  colback=blue!3, colframe=info, coltitle=white, fonttitle=\bfseries,
  title={\faCheckCircle\ #1}, breakable, arc=4pt
}
\newtcolorbox{adminbox}{
  colback=purple!5, colframe=admin, coltitle=white, fonttitle=\bfseries,
  title={\faCrown\ Solo Administrador}, breakable, arc=4pt
}

% ── Encabezado ───────────────────────────────────────────────────
\pagestyle{fancy}
\fancyhf{}
\fancyhead[L]{\textcolor{primario}{\textbf{MUYMUY Beauty Studio}}}
\fancyhead[R]{\textcolor{gris}{\small Manual de Administrador}}
\fancyfoot[C]{\textcolor{gris}{\small Página \thepage}}
\renewcommand{\headrulewidth}{0.5pt}
\renewcommand{\headrule}{\hbox to\headwidth{\color{primario}\leaders\hrule height \headrulewidth\hfill}}

\hypersetup{colorlinks=true, linkcolor=primario, urlcolor=primario}

% ════════════════════════════════════════════════════════════════
\begin{document}

% ── Portada ─────────────────────────────────────────────────────
\begin{titlepage}
  \centering
  \vspace*{4cm}
  {\color{primario}\rule{\linewidth}{1pt}}\\[1cm]
  {\color{oscuro}\Huge\bfseries MUYMUY Beauty Studio}\\[0.5cm]
  {\color{primario}\Large Manual de Usuario}\\[0.3cm]
  {\color{oscuro}\LARGE\bfseries ROL: ADMINISTRADOR}\\[1cm]
  {\color{primario}\rule{\linewidth}{1pt}}\\[2cm]
  {\color{gris}\large Versión 1.1 --- \today}\\[0.5cm]
  {\color{gris}\normalsize Sistema Administrativo Interno}
  \vfill
  {\color{oscuro}\small Documento confidencial --- Solo para uso interno}
\end{titlepage}
\nopagecolor

\tableofcontents
\newpage

% ════════════════════════════════════════════════════════════════
\section{Introducción}

Este manual está dirigido a los \textbf{Administradores} del sistema de MUYMUY Beauty Studio. El rol de administrador tiene acceso completo a todas las secciones del sistema, incluyendo las operativas (que también usa el empleado) y las de gestión y análisis.

\begin{adminbox}
El administrador tiene acceso a \textbf{todas} las secciones: Inicio, Agenda, Clientes, Asistencia, Venta Directa, Caja, Vacaciones, \textbf{Análisis, Administración, Inventario} y \textbf{Marketing}.
\end{adminbox}

\begin{notabox}
Para las secciones operativas (Agenda, Clientes, Asistencia, Venta Directa, Caja), el funcionamiento es idéntico al descrito en el \textbf{Manual de Empleado}. Este manual cubre las \textbf{funciones exclusivas del administrador} y complementa las que ya conoces.
\end{notabox}

\subsection{Inicio de Sesión}

\begin{pasobox}{Acceder al sistema}
\begin{enumerate}[leftmargin=*]
  \item Entra a la URL del sistema en el navegador.
  \item Ingresa tu correo y contraseña de administrador.
  \item El sistema te llevará al \textbf{Dashboard de Inicio}.
\end{enumerate}
\end{pasobox}

% ════════════════════════════════════════════════════════════════
\section{Secciones Operativas (igual que Empleado)}

Las siguientes secciones funcionan de la misma manera que en el Manual de Empleado, pero con privilegios adicionales:

\subsection{Agenda (Funciones Adicionales de Admin)}

Como administrador, en la Agenda verás dos botones adicionales en la esquina superior derecha:

\begin{description}[leftmargin=0pt]
  \item[\textbf{Bloquear Horario:}] Abre un modal para crear bloqueos manuales en la agenda. Puedes:
    \begin{itemize}
      \item Seleccionar la empleada.
      \item Elegir el rango de fechas y horas.
      \item Escribir el motivo del bloqueo.
      \item Crear bloqueos recurrentes (días específicos de la semana).
    \end{itemize}
  \item[\textbf{Desbloqueo Masivo:}] Permite eliminar múltiples bloqueos de una vez. Útil cuando cambia el horario de una empleada o hay días festivos que se cancelan.
\end{description}

\begin{alertabox}
Úsalo con cuidado. Un bloqueo mal puesto puede hacer que los clientes no puedan agendar con una empleada disponible.
\end{alertabox}

% ════════════════════════════════════════════════════════════════
\newpage
\section{Dashboard de Inicio}

El Dashboard es la pantalla principal para el administrador. Muestra un resumen en tiempo real del estado del negocio.

\subsection{Selectores de Tiempo y Sucursal}

En la parte superior del Dashboard puedes filtrar por:
\begin{itemize}[leftmargin=*]
  \item \textbf{Período:} Hoy, Esta Semana, Este Mes.
  \item \textbf{Sucursal:} Ver todas las sucursales o una en específico.
\end{itemize}

\subsection{Tarjetas de KPI (Indicadores Clave)}

El Dashboard muestra las siguientes métricas:

\begin{center}
\begin{tabular}{@{}lp{8cm}@{}}
\toprule
\textbf{Indicador} & \textbf{Descripción} \\
\midrule
Ingresos Totales & Suma de todos los tickets pagados en el período. \\
Citas & Número total de citas registradas. \\
Nuevos Clientes & Clientes que visitaron por primera vez. \\
Tasa de Asistencia & Porcentaje de citas finalizadas vs. total. \\
Ticket Promedio & Ingreso promedio por visita. \\
Tasa de Retención & Porcentaje de clientas que regresan. \\
Ocupación & Porcentaje de slots de agenda utilizados. \\
\bottomrule
\end{tabular}
\end{center}

\subsection{Gráficas del Dashboard}

El Dashboard incluye gráficas interactivas:
\begin{itemize}[leftmargin=*]
  \item \textbf{Ingresos por Sucursal:} Comparativa de ingresos, sueldos y comisiones entre sucursales.
  \item \textbf{Top 10 Empleadas:} Las profesionales con mayor generación de ingreso.
  \item \textbf{Tendencia de Servicios:} Servicios más solicitados por categoría a lo largo del tiempo.
  \item \textbf{Métodos de Pago:} Distribución de pagos por forma de pago.
\end{itemize}

\begin{notabox}
Los datos del Dashboard se actualizan automáticamente. Para forzar un refresco manual, hay un botón de actualizar en la esquina superior derecha.
\end{notabox}

% ════════════════════════════════════════════════════════════════
\newpage
\section{Análisis y Reportes}

La sección de Análisis es la herramienta de inteligencia de negocio del sistema. Está dividida en tres pestañas.

\subsection{Pestaña: Ventas y Cajas}

Contiene reportes financieros detallados:

\begin{itemize}[leftmargin=*]
  \item Ingresos por período (diario, semanal, mensual).
  \item Facturación neta (ingresos menos gastos de sueldos y comisiones).
  \item Análisis de caja (diferencias entre efectivo esperado y real).
  \item Reporte de métodos de pago.
  \item Facturación por tratamiento.
  \item Facturación estimada (citas agendadas pero no cobradas aún).
\end{itemize}

\subsection{Pestaña: Estadísticas y Reportes}

Indicadores de desempeño operativo:

\begin{itemize}[leftmargin=*]
  \item Nuevos clientes por período.
  \item Tasa de no-show (clientes que no asistieron).
  \item Análisis de asistencia del personal.
  \item Tendencia de citas agendadas vs. realizadas.
  \item Heatmap de afluencia por hora y día de la semana.
  \item Mix de servicios por familia.
  \item Ticket promedio por sucursal.
  \item Retención de clientes.
\end{itemize}

\subsection{Pestaña: Desempeño y Comisiones}

Reportes de rendimiento del equipo:

\begin{itemize}[leftmargin=*]
  \item Facturación individual por empleada.
  \item Cálculo de comisiones por ventas (basado en umbrales configurados).
  \item Evaluación de cumplimiento de hoja de desempeño.
  \item Ingresos por sucursal desglosados.
\end{itemize}

\subsection{Cómo Generar un Reporte}

\begin{pasobox}{Pasos para calcular un reporte}
\begin{enumerate}[leftmargin=*]
  \item Ve a \textbf{Análisis} y selecciona la pestaña correspondiente.
  \item En el selector, elige el \textbf{indicador} que deseas ver.
  \item Configura el \textbf{rango de fechas} (inicio y fin).
  \item Selecciona la \textbf{sucursal} (o ``Todas'').
  \item Elige el tipo de \textbf{desglose} (por sucursal, por empleada, por fecha, etc.).
  \item Selecciona el \textbf{ordenamiento} de los resultados.
  \item Haz clic en \textbf{Calcular}.
  \item \textit{[EN DESARROLLO]} Opcionalmente, haz clic en \textbf{Exportar a Excel} para descargar el reporte en formato CSV.
\end{enumerate}
\end{pasobox}

% ════════════════════════════════════════════════════════════════
\newpage
\section{Administración}

La sección de Administración es el panel de configuración del sistema.

\subsection{Gestión de Empleadas}

Desde aquí puedes administrar el equipo de trabajo:

\begin{pasobox}{Registrar una nueva empleada}
\begin{enumerate}[leftmargin=*]
  \item Ve a \textbf{Administración} y selecciona la pestaña de empleadas.
  \item Haz clic en \textbf{Nueva Empleada}.
  \item Llena los datos: nombre completo, sucursal asignada, fecha de contratación y sueldo diario.
  \item Guarda. La empleada aparecerá en la agenda de su sucursal.
\end{enumerate}
\end{pasobox}

También puedes:
\begin{itemize}[leftmargin=*]
  \item Editar datos de una empleada existente.
  \item Desactivar empleadas que ya no trabajan en el salón (sin borrar su historial).
\end{itemize}

\subsection{Gestión de Usuarios del Sistema}

\begin{pasobox}{Crear un usuario de acceso}
\begin{enumerate}[leftmargin=*]
  \item En \textbf{Administración}, ve a la sección de Usuarios.
  \item Haz clic en \textbf{Nuevo Usuario}.
  \item Ingresa el correo electrónico, nombre y rol (\textit{empleado} o \textit{admin}).
  \item El sistema enviará un correo de invitación para que establezcan su contraseña.
\end{enumerate}
\end{pasobox}

\begin{alertabox}
El rol de \textbf{superadmin} tiene los mismos permisos que admin pero no puede ser modificado desde el panel. Úsalo con extrema precaución.
\end{alertabox}

\subsection{Gestión de Sucursales}

Puedes registrar o editar las sucursales del negocio:
\begin{itemize}[leftmargin=*]
  \item Nombre de la sucursal.
  \item Dirección completa.
  \item Teléfono.
  \item RFC asociado.
  \item Número de cabinas disponibles.
\end{itemize}

\subsection{Gestión de Servicios}

Administra el catálogo de servicios del salón:
\begin{itemize}[leftmargin=*]
  \item Nombre del servicio.
  \item Familia o categoría (Esmaltado Permanente, Uñas Esculpidas, Spa, etc.).
  \item Precio.
  \item Duración en slots (cada slot = 30 minutos).
  \item Activar o desactivar servicios del catálogo.
\end{itemize}

\subsection{Configuración de Comisiones}

El sistema calcula comisiones automáticamente basado en umbrales de venta:

\begin{center}
\begin{tabular}{@{}ccc@{}}
\toprule
\textbf{Umbral de Ventas} & \textbf{\% Con Hoja} & \textbf{\% Sin Hoja} \\
\midrule
Hasta \$24,000 & 4.0\% & 2.0\% \\
\$24,001 -- \$30,000 & 4.5\% & 2.5\% \\
\$30,001 -- \$36,000 & 5.0\% & 3.0\% \\
... & ... & ... \\
Más de \$84,000 & 9.0\% & 7.0\% \\
\bottomrule
\end{tabular}
\end{center}

\textit{``Con Hoja''} significa que la empleada cumplió su evaluación de desempeño del mes. Puedes modificar estos porcentajes desde la sección de configuración de comisiones.

% ════════════════════════════════════════════════════════════════
\newpage
\section{Inventario}

La sección de Inventario permite gestionar los productos que se venden en el salón.

\subsection{Vista del Inventario}

El inventario muestra:
\begin{itemize}[leftmargin=*]
  \item Lista de todos los productos activos.
  \item Stock actual de cada producto.
  \item \textbf{Semáforo de stock:} Verde (10+), Amarillo (4--9), Rojo (0--3).
  \item Precio de venta y precio de costo.
\end{itemize}

\subsection{Gestionar Productos}

\begin{pasobox}{Agregar un producto nuevo}
\begin{enumerate}[leftmargin=*]
  \item Ve a \textbf{Inventario}.
  \item Haz clic en \textbf{Nuevo Producto}.
  \item Ingresa: nombre, precio de venta, precio de costo y stock inicial.
  \item Guarda. El producto estará disponible para venta directa y en los tickets de cita.
\end{enumerate}
\end{pasobox}

\begin{pasobox}{Actualizar stock}
\begin{enumerate}[leftmargin=*]
  \item Encuentra el producto en la lista.
  \item Haz clic en \textbf{Editar}.
  \item Modifica la cantidad de stock.
  \item Guarda los cambios.
\end{enumerate}
\end{pasobox}

\begin{notabox}
El stock se descuenta automáticamente cada vez que se vende un producto en una cita o venta directa.
\end{notabox}

% ════════════════════════════════════════════════════════════════
\newpage
\section{Marketing [EN DESARROLLO]}

La sección de Marketing permitirá gestionar campañas de comunicación con los clientes. Actualmente se encuentra en fase de desarrollo.

\subsection{Campañas de WhatsApp [PRÓXIMAMENTE]}

Puedes enviar mensajes masivos o segmentados a tus clientes:

\begin{pasobox}{Crear una campaña}
\begin{enumerate}[leftmargin=*]
  \item Ve a \textbf{Marketing}.
  \item Haz clic en \textbf{Nueva Campaña}.
  \item Escribe el mensaje que deseas enviar.
  \item Selecciona el segmento de clientes:
    \begin{itemize}
      \item Todos los clientes.
      \item Clientes que no han visitado en X días.
      \item Clientes de una sucursal específica.
      \item Clientes con cumpleaños en el mes.
    \end{itemize}
  \item Revisa la lista de destinatarios.
  \item Confirma el envío.
\end{enumerate}
\end{pasobox}

\subsection{Configuración de Mensajes Automáticos [PRÓXIMAMENTE]}

Se podrán configurar mensajes automáticos que el sistema envía sin intervención manual:
\begin{itemize}[leftmargin=*]
  \item \textbf{Confirmación de cita:} Envío al agendar.
  \item \textbf{Recordatorio 24h:} Envío el día anterior.
  \item \textbf{Felicitación de cumpleaños:} Envío el día del cumpleaños.
\end{itemize}

% ════════════════════════════════════════════════════════════════
\newpage
\section{Gestión de Vacaciones}

Como administrador, eres el responsable de aprobar o rechazar las solicitudes de tiempo libre del equipo.

\subsection{Revisar Solicitudes}

En la sección \textbf{Vacaciones} verás un listado de todas las solicitudes de tu sucursal. Las solicitudes pendientes aparecerán con botones de acción rápida.

\begin{pasobox}{Aprobar una solicitud}
\begin{enumerate}[leftmargin=*]
  \item Busca la solicitud con estado \textbf{Pendiente}.
  \item Revisa las fechas y las notas de la empleada.
  \item Haz clic en \textbf{Aprobar}.
  \item Opcionalmente, escribe una nota para la empleada.
  \item Confirma. El sistema \textbf{creará automáticamente los bloqueos} en la agenda de la empleada para esas fechas.
\end{enumerate}
\end{pasobox}

\begin{pasobox}{Rechazar una solicitud}
\begin{enumerate}[leftmargin=*]
  \item Haz clic en \textbf{Rechazar} en la solicitud correspondiente.
  \item \textbf{Importante:} Escribe el motivo del rechazo en el campo de notas para que la empleada sepa la razón.
  \item Confirma el rechazo.
\end{enumerate}
\end{pasobox}

\subsection{Automatización con la Agenda}

El módulo de vacaciones está integrado directamente con la agenda:
\begin{itemize}[leftmargin=*]
  \item Al aprobar una solicitud, se generan bloqueos de jornada completa (08:00 a 21:00) en todos los días del rango.
  \item Si se aprueba una \textbf{Extensión}, se agregarán nuevos bloqueos para los días adicionales.
  \item Esto garantiza que no se agenden citas accidentalmente mientras la profesional está fuera.
\end{itemize}

\begin{notabox}
Los administradores actúan únicamente como aprobadores. El botón de ``Nueva Solicitud'' está oculto para el rol admin para evitar confusiones.
\end{notabox}

% ════════════════════════════════════════════════════════════════
\newpage
\section{Evaluaciones de Desempeño (Hoja)}

El sistema permite registrar si cada empleada ``cumplió su hoja'' mensualmente. Esto afecta directamente el porcentaje de comisión que recibe.

\begin{pasobox}{Registrar evaluación mensual}
\begin{enumerate}[leftmargin=*]
  \item Ve a \textbf{Administración} y busca la sección de Evaluaciones.
  \item Selecciona el mes y año a evaluar.
  \item Para cada empleada, indica si \textbf{cumplió} o \textbf{no cumplió} su hoja de desempeño.
  \item Guarda. El sistema usará estos datos para calcular las comisiones correctamente en los reportes.
\end{enumerate}
\end{pasobox}

% ════════════════════════════════════════════════════════════════
\newpage
\section{Preguntas Frecuentes del Administrador}

\begin{description}[leftmargin=0pt]
  \item[\textbf{¿Cómo reseteo la contraseña de una empleada?}]
    Ve a Administración \textrightarrow\ Usuarios, selecciona a la empleada y usa la opción ``Restablecer contraseña''. Se enviará un correo de reseteo.

  \item[\textbf{¿Cómo veo los reportes de todas las sucursales a la vez?}]
    En la sección de Análisis, elige ``Todas las sucursales'' en el filtro de sucursal antes de calcular.

  \item[\textbf{¿Se pueden recuperar datos borrados?}]
    Los empleados y clientes desactivados conservan su historial. Los tickets y citas no se eliminan, solo cambian de estado. Contacta soporte técnico si necesitas recuperar datos específicos.

  \item[\textbf{¿Con qué frecuencia se actualizan los datos del Dashboard?}]
    El Dashboard tiene un caché de 3 minutos. Los datos del análisis se calculan en tiempo real al hacer clic en ``Calcular''.

  \item[\textbf{¿Cómo cambio el precio de un servicio?}]
    Ve a Administración \textrightarrow\ Servicios, edita el servicio y actualiza el precio. Afectará las nuevas citas; las ya creadas conservan el precio original.

  \item[\textbf{¿Puedo tener múltiples administradores?}]
    Sí. Puedes crear varios usuarios con rol ``admin''. Todos tendrán acceso completo al sistema.
\end{description}

% ════════════════════════════════════════════════════════════════
\vspace{2cm}
\begin{center}
  \textcolor{gris}{\rule{0.5\linewidth}{0.4pt}}\\[0.3cm]
  \textcolor{gris}{\small MUYMUY Beauty Studio --- Manual de Administrador v1.1}\\
  \textcolor{primario}{\small Documento de uso interno y confidencial}
\end{center}

\end{document}
